% ==============================================================================
% Fichier : front/03b_labs.tex : SGSI (IS498)
% ==============================================================================

\section*{Les dix labs de NDAP SUKULU}

\begin{tableaucours}{L{2.4cm}L{8.4cm}L{1.1cm}}{%
    \tableauHeader{Lab} &
    \tableauHeader{Sujet} &
    \tableauHeader{Ch.}
}
SGSI-LAB-01 & Inventaire et classification des actifs & 1 \\
SGSI-LAB-02 & Rédiger une politique de sécurité & 5 \\
SGSI-LAB-03 & Analyse d'écart ISO 27001 (gap analysis) & 8 \\
SGSI-LAB-04 & Appréciation des risques (EBIOS Risk Manager) & 12 \\
SGSI-LAB-05 & Matrice et traitement des risques & 13 \\
SGSI-LAB-06 & Construire une déclaration d'applicabilité & 14 \\
SGSI-LAB-07 & Contrôle d'accès et moindre privilège (RBAC) & 16 \\
SGSI-LAB-08 & Plan de continuité et sauvegarde 3-2-1 & 17 \\
SGSI-LAB-09 & Tableau de bord d'indicateurs de sécurité & 18 \\
SGSI-LAB-10 & Audit interne et non-conformités & 19 \\
\end{tableaucours}

\medskip
Chaque lab dure environ une heure et se corrige automatiquement. Les pièces
se trouvent dans le dossier \texttt{preuves} (lecture seule) et une copie de
travail dans \texttt{travail}.
